% !TeX root = ../main.tex
\chapter{Introduction to Line-less Mobile Assembly Systems}
\label{ch:introduction}

% =====================================================================
\section{Line-less Mobile Assembly Systems and Their Implications for
  Motion Planning}
\label{sec:lmas}

Manufacturing operates under shortening product life cycles, growing
variant diversity, and demand that fluctuates over the life of a
product; a CIRP-level review of product-variety management documents
this trend across industries and quantifies the cost it imposes on
assembly~\cite{ElMaraghy2013Variety}. Rigid, station-based assembly
lines are dimensioned for a fixed product mix and a stable takt time.
When the product or the required volume changes, layout, material flow,
and station balancing have to be reconfigured, which is slow and
capital-intensive~\cite{Huettemann2016LMAS}.

Manufacturing research has answered this tension with a family of
paradigms that trade dedicated throughput for adaptability. Flexible
Manufacturing Systems provide a fixed set of capabilities that can be
re-combined for a pre-defined product family, whereas Reconfigurable
Manufacturing Systems are designed so that capacity and functionality
themselves can be changed by adding, removing, or rearranging
modules~\cite{Koren1999RMS}. ElMaraghy contrasts the two paradigms and
places them on a common axis of changeability, with flexibility bought
in advance in the former case and acquired on demand in the
latter~\cite{ElMaraghy2005FRMS}. In assembly specifically, this line of
work has produced modular station concepts in which workstations are no
longer traversed in a fixed sequence~\cite{Kern2016Modular}.

The \emph{Line-less Mobile Assembly System} (LMAS) pushes this
development to its limit by removing the fixed station grid
altogether~\cite{Huettemann2019LMAS}. Instead of binding products,
tools, and workers to a sequence of static stations, an LMAS models all
production resources---mobile robots, mobile manipulators, workpiece
carriers, and human workers---as \emph{mobile, freely relocatable
entities} on a shared shop floor, whose assignment, sequence, and
location are planned at run time~\cite{Buckhorst2019Planning}. The
assembly ``line'' thus becomes a reconfigurable constellation of
resources matched to the current order, and control of that
constellation is decentralised so that the system can absorb product
changes, demand fluctuations, and disturbances without a physical
relayout~\cite{Buckhorst2022LMAS}.

A central enabler of the LMAS concept is the \emph{mobile manipulator}:
a robotic arm mounted on a mobile base
(\cref{fig:mobile-manipulator}). The mobile base extends the reachable
workspace far beyond the arm's own envelope and lets the robot
reposition itself between and within tasks, while the arm provides the
dexterity required for the assembly operation itself. Surveys of the
field identify exactly this combination (one resource servicing many
workstations, covering workpieces larger than a fixed-base arm could
reach) as the motivation for wheeled mobile
manipulation~\cite{Sandakalum2022Survey,Thakar2023MobileManip}.

\begin{figure}[htbp]
  \centering
  \begin{tikzpicture}[
      font=\footnotesize,
      lbl/.style={font=\scriptsize, align=center},
      dim/.style={{Latex[length=1.5mm]}-{Latex[length=1.5mm]}, thin},
      link/.style={line width=1.1pt, line cap=round},
      jt/.style={circle, draw, fill=white, inner sep=0pt, minimum size=2.0mm},
      ax/.style={-{Latex[length=1.6mm]}, thick},
    ]
    % =================================================================
    % (a) top view: planar base pose and the no-slip direction
    % =================================================================
    \begin{scope}[shift={(0,0)}]
      \draw[ax] (0,0) -- (0.85,0) node[lbl, right=0.5mm] {$x_W$};
      \draw[ax] (0,0) -- (0,0.85) node[lbl, above=0.5mm] {$y_W$};
      \node[lbl, below left=0mm] at (0,0) {$W$};

      \begin{scope}[shift={(2.6,1.75)}, rotate=25]
        \draw[fill=gray!12] (-0.92,-0.60) rectangle (0.92,0.60);
        \draw[fill=gray!40] (-0.24,-0.60) rectangle (0.24,-0.76);
        \draw[fill=gray!40] (-0.24, 0.60) rectangle (0.24, 0.76);
        \draw[ax] (0,0) -- (1.45,0) node[lbl, above right=-0.4mm] {$x_B$};
        \draw[ax] (0,0) -- (0,1.20) node[lbl, above=0.5mm] {$y_B$};
        % blocked lateral direction
        \draw[dashed, gray, {Latex[length=1.4mm]}-{Latex[length=1.4mm]}]
          (0,-1.35) -- (0,-0.80);
        \draw[thick] (-0.13,-1.20) -- (0.13,-1.46);
        \draw[thick] (-0.13,-1.46) -- (0.13,-1.20);
      \end{scope}
      \node[lbl, anchor=west] at (3.55,0.55) {no side-slip};

      \draw[dim] (0,-0.50) -- (2.6,-0.50) node[lbl, midway, below] {$x_b$};
      \draw[dim] (-0.50,0) -- (-0.50,1.75) node[lbl, midway, left] {$y_b$};
      \draw[dashed, gray] (2.6,1.75) -- (2.6,-0.50);
      \draw[dashed, gray] (2.6,1.75) -- (-0.50,1.75);
      \draw[dashed, gray] (2.6,1.75) -- ++(1.75,0);
      \draw[-{Latex[length=1.4mm]}] (3.85,1.75) arc (0:25:1.25);
      \node[lbl] at (4.05,2.12) {$\theta_b$};

      \node[lbl, anchor=north] at (1.9,-1.10) {(a) base pose in the plane};
    \end{scope}

    % =================================================================
    % (b) side view: base plus six-joint arm
    % =================================================================
    \begin{scope}[shift={(6.9,-0.60)}]
      \draw[thick] (-0.4,-0.25) -- (5.0,-0.25);
      \foreach \x in {-0.25,0.10,...,4.90}
        \draw[gray, thin] (\x,-0.25) -- ++(-0.16,-0.16);

      \draw[fill=gray!12] (0.0,0.12) rectangle (2.3,0.95);
      \draw[fill=gray!40] (0.5,0.12) circle (0.37);
      \draw[fill=gray!40] (1.8,0.12) circle (0.37);
      \node[lbl] at (1.15,0.60) {base};

      \coordinate (a0) at (1.75,0.95);
      \coordinate (a1) at (1.75,1.70);
      \coordinate (a2) at (3.00,2.45);
      \coordinate (a3) at (4.05,1.75);
      \coordinate (a4) at (4.60,2.08);
      \draw[link] (a0) -- (a1) -- (a2) -- (a3) -- (a4);
      \node[jt] at (a0) {}; \node[jt] at (a1) {};
      \node[jt] at (a2) {}; \node[jt] at (a3) {};
      \node[lbl, anchor=south] at (2.35,2.70)
        {UR10e: six revolute\\joints $q_1,\dots,q_6$};

      \draw[ax] (a4) -- ++(0.60,0.25);
      \draw[ax] (a4) -- ++(-0.25,0.60);
      \node[lbl, anchor=north west] at (4.62,1.90) {\ee{}\\pose in $SE(3)$};

      \node[lbl, anchor=north] at (2.3,-0.75) {(b) arm mounted on the base};
    \end{scope}
  \end{tikzpicture}
  \caption[Mobile manipulator: base and arm degrees of freedom]{%
    Wheeled mobile manipulator of the type studied here. (a)~The base
    contributes the planar pose $(x_b,y_b,\theta_b)$; its wheels cannot
    slide sideways, which constrains the velocities it can produce but
    not the poses it can reach. (b)~The arm contributes six revolute
    joints, and the task is the pose of the \ee{} frame in $SE(3)$.
    Counting these groups against the task is what makes the system
    redundant (\cref{eq:grubler,eq:redundancy}).}
  \label{fig:mobile-manipulator}
\end{figure}

Adding a mobile base to the arm changes the motion-planning problem in
three ways. In a classical line, an industrial arm operates in an
engineered, static cell: the part arrives at a known pose, the
surroundings are fixed, and the arm executes largely pre-programmed
trajectories. In an LMAS, three of the assumptions behind that mode of
operation are weakened.

\paragraph{Large, base-dependent workspace.}
The target \ee{} pose may lie outside the reach of the arm from the
current base pose, so base and arm have to \emph{cooperate} to
accomplish a single task. Both mobile-manipulation surveys treat the
choice of a base pose that renders the task reachable and
well-conditioned as a planning problem in its own
right~\cite{Sandakalum2022Survey,Thakar2023MobileManip}.

\paragraph{Shared workspace with moving resources.}
Because resources are relocatable, parts of the environment move
relative to one another during operation. How much of the environment is
known in advance depends on the plant: a shop floor with tracked
resources and a maintained world model is partially predictable, whereas
human co-workers and manually moved carriers are not. The consequence
for planning is therefore not that offline computation is impossible,
but that a trajectory computed once for an assumed scene has to be
revised while it is executed, which is the case that motivates online
re-planning~\cite{Finean2023ContextAware}.

\paragraph{Kinematic redundancy.}
The base adds degrees of freedom (DOF) to a system that already has as
many as the task requires. The arm is a spatial open chain with
$N = 7$ links (including its mounting link) connected by $J = 6$
revolute joints, each contributing $f_i = 1$ freedom. Gr\"ubler's
formula for a spatial mechanism ($m = 6$) gives its
mobility~\cite{Lynch2017ModernRobotics}
\begin{equation}
  \label{eq:grubler}
  \mathrm{dof} = m\,(N - 1 - J) + \sum_{i=1}^{J} f_i
               = 6\,(7 - 1 - 6) + 6 = 6 .
\end{equation}
The base contributes a planar pose, so the whole-body configuration is
$q = (x_b, y_b, \theta_b, q_1, \dots, q_6)$ and the configuration space
is nine-dimensional, while the \ee{}-pose task lives in the
six-dimensional $SE(3)$:
\begin{equation}
  \label{eq:redundancy}
  \dim \mathcal{Q} - \dim \mathcal{T} = 9 - 6 = 3 .
\end{equation}
Three configuration-space directions therefore leave the \ee{} pose
unchanged. The picture differs at velocity level, because a
differential-drive base cannot slide sideways: the rolling-without-slip
condition $\dot{x}_b \sin\theta_b - \dot{y}_b \cos\theta_b = 0$ is a
non-holonomic constraint, so the base offers only two instantaneous
inputs (a forward and a turning rate) while still reaching all three
planar coordinates over time~\cite{Lynch2017ModernRobotics,%
Siciliano2009Robotics}. The instantaneous input vector is therefore
eight-dimensional against a six-dimensional task twist, leaving two
degrees of instantaneous redundancy. This redundancy is an opportunity:
the same \ee{} motion can be produced by infinitely many base--arm
combinations, and the surplus can be spent on reach, on
manipulability~\cite{Yoshikawa1985Manipulability}, or on obstacle
avoidance~\cite{Dietrich2015NullSpace}. It is an opportunity only if the
planner reasons about base and arm \emph{jointly}; a planner that fixes
the base first spends the redundancy before it can be used.

Together, these three conditions call for motion planning that is
reactive, that spans the \emph{combined} configuration space of base and
arm, and that resolves the resulting redundancy deliberately rather than
as a by-product. Whole-body motion planning, introduced next, addresses
exactly this problem.

% =====================================================================
\section{The Need for Whole-Body Motion Planning: From Sequential to
  Unified Approaches}
\label{sec:need-wholebody}

The most common way to command a mobile manipulator is to treat base and
arm as two independent subsystems and to plan for them
\emph{sequentially}: a navigation stack such as
Nav2~\cite{Macenski2020Nav2} first drives the base to a suitable
stand-off pose; an arm-motion framework such as
MoveIt~\cite{moveit2} then solves the inverse kinematics for the target
\ee{} pose from that now-fixed base pose and plans a joint-space
trajectory to the resulting configuration, typically with a
sampling-based planner~\cite{Sucan2012OMPL}. This decoupled scheme is
modular and, above all, built from components that are mature and
widely deployed in industry, which is the practical reason for its
prevalence; both mobile-manipulation surveys report it as the dominant
architecture in fielded
systems~\cite{Sandakalum2022Survey,Thakar2023MobileManip}.

Decoupling has three consequences that become severe in the LMAS
setting; \cref{fig:decoupled-vs-wholebody} contrasts the two schemes.

\paragraph{Loss of reach and of solution quality.}
Choosing the base pose \emph{before} planning the arm discards the
coupling between the two subsystems. A base pose that looks reasonable
geometrically may leave the arm near a singularity or a joint limit, or
may make the \ee{} target unreachable, even though a jointly chosen
base--arm configuration
exists~\cite{Sandakalum2022Survey,Thakar2023MobileManip}. Both surveys
describe base placement and arm planning as coupled sub-problems, and
the holistic controller of Haviland et~al.\ is motivated by the same
observation~\cite{Haviland2022Holistic}.

\paragraph{Stop-and-go behaviour.}
Sequential execution---drive, stop, then manipulate---prevents base and
arm from moving \emph{simultaneously}. Tasks in which the \ee{} has to
follow a continuous path longer than the arm's reach, which is a common
LMAS case, cannot be executed smoothly this
way~\cite{BurgessLimerick2024ReactiveBase,Pankert2020PerceptiveMPC}.
Reported gains of on-the-move manipulation over the stop-and-go pipeline
motivate both the reactive base controller of Burgess-Limerick
et~al.\ and the continuous mobile manipulation of Pankert and Hutter.

\paragraph{Cost of reacting to change.}
A plan computed once, for a fixed base pose, does not by itself account
for base drift, moving obstacles, or an updated target. This is a
question of degree rather than of possibility: the pipeline can be
re-triggered, and planners that predict how the scene will evolve and
repair the trajectory online exist~\cite{Finean2023ContextAware}. What a
decoupled pipeline cannot do cheaply is revise the \emph{base} decision,
because that decision sits upstream of the arm plan and invalidates it.

\begin{figure}[htbp]
  \centering
  \begin{tikzpicture}[
      font=\footnotesize,
      box/.style={draw, rounded corners=1pt, align=center,
                  inner sep=3.5pt, minimum height=8mm, minimum width=22mm},
      flow/.style={-{Latex[length=2mm]}, thick},
      lbl/.style={font=\scriptsize, align=center},
      lane/.style={draw, dashed, rounded corners=2pt, inner sep=6pt},
    ]
    % ---------- decoupled lane --------------------------------------
    \node[box] (d1) at (0,0)    {plan base\\path};
    \node[box] (d2) at (2.9,0)  {drive base,\\arm parked};
    \node[box] (d3) at (5.8,0)  {stop,\\solve IK};
    \node[box] (d4) at (8.7,0)  {plan and\\execute arm};
    \draw[flow] (d1) -- (d2); \draw[flow] (d2) -- (d3);
    \draw[flow] (d3) -- (d4);
    \node[lane, fit=(d1)(d4),
          label={[lbl]above:\textbf{Decoupled:} base and arm move in turn}] {};
    \node[lbl, anchor=north] at (5.8,-0.75)
      {a new base pose invalidates the arm plan};

    % ---------- whole-body lane -------------------------------------
    \node[box, minimum width=34mm] (w1) at (1.1,-3.2)
      {\ee{} reference\\in task space};
    \node[box, minimum width=40mm] (w2) at (5.9,-3.2)
      {one optimisation over\\$(x_b,y_b,\theta_b,q_1,\dots,q_6)$};
    \node[box, minimum width=22mm] (w3a) at (9.9,-2.65) {base twist};
    \node[box, minimum width=22mm] (w3b) at (9.9,-3.75) {arm rates};
    \draw[flow] (w1) -- (w2);
    \draw[flow] (w2.east) -- (w3a.west);
    \draw[flow] (w2.east) -- (w3b.west);
    \draw[flow] ([yshift=-2mm]w2.south) -- ++(0,-0.7) -| ([yshift=-2mm]w1.south)
      node[lbl, pos=0.25, below] {state feedback, re-solve each step};
    \node[lane, fit=(w1)(w2)(w3a)(w3b),
          label={[lbl]above:\textbf{Whole-body:} base and arm move together}] {};
  \end{tikzpicture}
  \caption[Decoupled versus whole-body command of a mobile
    manipulator]{%
    Decoupled versus whole-body command of a mobile manipulator. The
    decoupled pipeline fixes the base pose upstream of the arm plan, so
    the two subsystems act in sequence. A whole-body scheme optimises
    over the joint configuration, emits base and arm commands from the
    same solution, and closes the loop on the measured state.}
  \label{fig:decoupled-vs-wholebody}
\end{figure}

\emph{Whole-body motion planning} removes this boundary by treating
mobile base and arm as a \emph{single} redundant kinematic (or dynamic)
system and planning their motion in one
optimisation~\cite{Lynch2017ModernRobotics,Haviland2022Holistic}. The
task is expressed directly in task space and the planner distributes the
required motion between base and arm according to a chosen cost
criterion, resolving the redundancy of \cref{eq:redundancy} explicitly
rather than by construction. The formulation generalises redundancy
resolution for kinematically redundant
manipulators~\cite{Siciliano2009Robotics} and null-space projection for
a fixed base~\cite{Dietrich2015NullSpace} to base and arm together.
This addresses the three points above: the full combined workspace
becomes available, base and arm move simultaneously along one
coordinated trajectory, and the base decision is revisited at every
re-solve instead of being frozen upstream. How well the last property
holds in practice depends on how fast the optimisation can be re-solved
relative to the rate at which the scene changes; this work quantifies
that trade-off on hardware in \cref{ch:hardware}.

Formulating whole-body planning as a receding-horizon optimal-control
problem, that is, as \emph{Model Predictive Control}
(\mpc{})~\cite{Rawlings2017MPC}, suits mobile manipulation for two
reasons. First, the constraints that make the problem hard---joint and
velocity limits, self-collision, and the non-holonomic condition of a
wheeled base---are inequality constraints on states and inputs, which is
the form an \mpc{} formulation handles
natively~\cite{Mayne2000Constrained}. Second, the solver re-plans from
the measured state at every step, so the same mechanism that produces
the plan also provides the feedback. Whole-body \mpc{} on this basis has
been demonstrated on a dynamically stable mobile
manipulator~\cite{Minniti2019WholeBodyMPC} and for continuous mobile
manipulation from onboard perception~\cite{Pankert2020PerceptiveMPC}.
The open-source \ocs{}~(Optimal Control for Switched Systems)
framework~\cite{ocs2Toolbox} provides an implementation of this approach
and has been used for whole-body loco-manipulation on legged and wheeled
platforms~\cite{Sleiman2021UnifiedMPC}.

% =====================================================================
\section{Objectives and Structure of This Work}
\label{sec:objectives}

\subsection{Objective}
\label{sec:obj-goal}
The goal of this work is to \emph{integrate and validate} whole-body
\mpc{}, based on the \ocs{} framework, on the \hero{} mobile
manipulator (a differential-drive mobile base carrying a Universal
Robots UR10e arm~\cite{UR10e}), and to characterise the resulting
\ee{}-pose tracking behaviour under realistic conditions. In contrast to
studies that remain in simulation, the emphasis is on a
\emph{deployable} system: the \mpc{} runs in a distributed, real-time
software architecture on the physical robot, and its performance is
assessed on hardware.

\subsection{Requirements on the Motion Planner}
\label{sec:requirements}
The discussion above yields the criteria that any candidate planner has
to satisfy on \hero{}. They are stated here explicitly, so that the
comparison in \cref{sec:sota-selection} has a fixed reference and does
not rest on an implicit notion of suitability. \Cref{tab:requirements}
lists them.

\begin{table}[htbp]
  \centering
  \caption[Requirements on the motion planner for \hero{}]{%
    Requirements on the motion planner for \hero{}, with the argument
    each one follows from. \Cref{tab:sota-comparison} scores the
    candidate approaches against exactly these criteria.}
  \label{tab:requirements}
  \begin{tabularx}{\linewidth}{@{}llX@{}}
    \toprule
    ID & Requirement & Origin \\
    \midrule
    R1 & Coordinated base--arm motion
       & Reach and stop-and-go arguments, \cref{sec:need-wholebody} \\
    R2 & Explicit constraint handling
       & Joint, velocity, non-holonomic and self-collision limits,
         \cref{eq:redundancy} and \cref{sec:need-wholebody} \\
    R3 & Deliberate redundancy resolution
       & Three surplus configuration directions, \cref{eq:redundancy} \\
    R4 & Online re-planning from the measured state
       & Moving resources in a shared workspace, \cref{sec:lmas} \\
    R5 & Real-time feasibility on the on-board computer
       & Deployment on \hero{} rather than in simulation,
         \cref{sec:obj-goal} \\
    R6 & ROS~2 interfaces
       & Integration with the existing \hero{} software stack,
         \cref{ch:integration} \\
    R7 & Transparency and analysability
       & Behaviour has to be diagnosable on hardware,
         \cref{ch:hardware} \\
    \bottomrule
  \end{tabularx}
\end{table}

\subsection{Scope and Contributions}
\label{sec:obj-contrib}
The concrete contributions of this thesis are:
\begin{enumerate}
  \item the selection of a suitable whole-body motion-planning approach
        for the \hero{} platform, following a comparison of candidate
        frameworks against \cref{tab:requirements} (\cref{ch:sota});
  \item the integration of \ocs{} whole-body \mpc{} into the \hero{}
        software stack, including a distributed host/on-board
        architecture for real-time operation, the whole-body model and
        its parameterisation, and the hardware and \mbox{ROS-control}
        interfaces (\cref{ch:integration});
  \item a simulation-based validation of the closed loop that isolates
        the geometry, timing, and data pipeline before hardware
        deployment, and that measures the residual the optimiser leaves
        on its own model (\cref{ch:simulation}); and
  \item the deployment of the system on the physical \hero{} platform,
        together with a discussion of the deployment challenges
        encountered and preliminary \ee{}-tracking results
        (\cref{ch:hardware}).
\end{enumerate}

\subsection{Structure}
\label{sec:obj-structure}
\Cref{ch:sota} reviews the state of the art in whole-body motion
planning for mobile manipulators, classifies the main approaches,
presents \ocs{} and alternative frameworks, and justifies the selection
made here. \Cref{ch:integration} describes the \hero{} platform and the
integration of \ocs{} into it: the distributed architecture, the
whole-body model and its parameterisation, and the interfaces to the
low-level control layer. \Cref{ch:simulation} presents the
simulation-based validation using a kinematic closed-loop (dummy MRT)
setup. \Cref{ch:hardware} reports on the hardware deployment, the
challenges and their mitigations, and preliminary \ee{}-tracking
results. \Cref{ch:conclusion} summarises the work, reflects critically
on its limitations, and gives recommendations for resolving the
remaining issues.
